\section{Closest Methods and the Remaining Question}
\label{sec:narrative-related}

Dynamic state estimation supplies the state and uncertainty machinery used to predict a PMU history. Unknown-input and DAE estimators address incomplete models and exogenous inputs \cite{gillijns2007input,nugroho2022dae,katanic2024recursive}. Their task is not automatically an event-family/source decision. A recovered bus-voltage field can remain compatible with competing physical interventions.

Model-based source inference is the closest methodological precedent. Phasor angles support outage hypotheses \cite{tate2008outage}; uncertain pre-event conditions can be marginalized with outage identity \cite{zhao2014outages}; and sequential likelihoods detect line changes under stochastic operation \cite{chen2016quickest,rovatsos2017outage}. Anguluri \emph{et al.} jointly recover sparse forcing locations, inputs and an unknown initial state in linear systems using group LASSO \cite{anguluri2023group}. That is a direct precedent for physical source recovery, although forced inputs and linear dynamics differ from the switched load/line episode task here. It is not an executed C01 comparator. Transient dynamics, uncertain initial states and unknown onset are therefore established ingredients, not contributions of this paper.

Warner \emph{et al.} provide a particularly close precedent \cite{warner2023highfidelity}. Their simulated voltage/current fault trajectories are compared with a steady-state reference over a moving onset window; Algorithm~1 stops and returns candidate lines. Load changes are an adverse-event test rather than a load-source output. Their Sections~II and IV target short-circuit identification before clearing, whereas our restricted task follows load changes and line openings through restoration and later interventions. Sparse sensing, a physical trajectory library and retrospective onset scanning are shared ideas. We have not run their procedure on the present data and make no numerical superiority claim over it.

Sequential diagnosis also has a stronger statistical literature than a rolling-window label suggests. The Adaptive Matrix CuSum addresses misidentification caused by pre-change observations and establishes delay/error results under its specified observation laws \cite{warner2024matrix}. R11's pruned conditional histories, colored measurement state and plug-in startup estimates do not inherit those guarantees. A ten-second window or the array $M$ is not a replacement for that theory.

Learned PMU diagnosis provides a different comparator. Sparse classifiers and knowledge-driven representations identify represented events or assets \cite{pandey2020realtime,li2024knowledge,yildiz2024few}. The historical SGSMA baseline supplies two frozen ExtraTrees adaptations here. We distinguish interpolation at represented sources, transfer of that frozen adapter, and true held-source generalization. The present causal catalog is known; its results do not establish the last of those claims.

Observability and diagnosability explain why neither a physical nor learned method is guaranteed to localize every visible event. Functional observability concerns the requested state function \cite{fernando2010functional,montanari2022functional}; descriptor-model diagnosability and nuisance-profiled tests concern competing explanations \cite{eriksson2013diagnosability,scharf1994matched}. Placement-dependent source ambiguity is also established in PMU fault/event work \cite{jamei2020limits,conte2023fault,yildiz2025sparse}. Our state/source distinction applies these ideas; it is not their first discovery.

Multiple-model filtering already manages conditional state/mode hypotheses \cite{blom1988imm}; R11 is a pruned history bank, not a new IMM derivation. Stochastic swing/kernel models estimate latent states in Ensaf and Barati \cite{ensaf2023kernel}, and physics-derived prior covariances support sparse-state forecasting \cite{tartakovsky2023gp}. Joint DAE-aware generator/network parameter inference is studied separately \cite{albustami2026dae}. R11 uses a finite bank of constant network hypotheses rather than inferring the continuous parameter posterior. A Gaussian process or neural architecture would require separate validation.

\begin{table*}[t]\centering\small
\caption{Closest inference contracts and the difference tested here. This is a methodological comparison; only Chen and the frozen ExtraTrees adapters are executed comparators in C01. Absence from a paper's evaluated task is not a claim of algorithmic impossibility.}\label{tab:literature-capability}
\begin{tabularx}{\textwidth}{@{}lXXX@{}}\toprule
Route & Physical/statistical mechanism & Published target & Difference or role here\\\midrule
Chen / Rovatsos & Increment statistics, network models and sequential tests & Line outages; Rovatsos includes transient phases & Chen is executed on persistent-line first stops\\
Warner high-fidelity & Simulated V/I paths, moving onset window, distance evidence & Short-circuit line candidates & Nearest trajectory precedent; no common-data run here\\
Adaptive Matrix CuSum & Pairwise sequential diagnosis and error/delay analysis & One unknown change among distributional alternatives & No corresponding optimality guarantee is claimed here\\
DAE / stochastic-state estimation & Incomplete DAE observers or physical uncertainty models & States, unknown inputs or network parameters & Components/adjacent tasks, not episode attribution scores\\
Knowledge-driven learned diagnosis & Spatiotemporal representation with topology/type constraints & Multi-family disturbance location & Distinct from the physical likelihood; ExtraTrees is our executed learned reference\\
Physical history bank here & Estimated-past-state paths, active/cause masses, fixed R/X mixture & Successive known-catalog load/line interventions & Test whether recovery-aware attribution retains timely sources without excess alarms\\\bottomrule
\end{tabularx}\end{table*}

The remaining question is operational and falsifiable: can a state-conditioned bank distinguish a new intervention from a preceding intervention's recovery while retaining timely family/source decisions under fixed network uncertainty? Our contribution is the specified recovery-aware construction and its bounded evidence, not a claim to have invented model-based diagnosis. We use ``physics-informed'' because physical equations define candidate trajectories and likelihood inputs. Learned discrepancy and the unrestricted joint state/event/integrity estimator remain separate proposals.
